\documentclass[11pt,a4paper]{article}
\usepackage[T1]{fontenc}
\usepackage[utf8]{inputenc}
\usepackage{lmodern,amsmath,amssymb}
\usepackage[margin=25mm]{geometry}
\usepackage{longtable,booktabs,array,calc}
\usepackage[hidelinks]{hyperref}
\hypersetup{pdftitle={The confinement scale in three bands, with published numbers, and the verification restate},pdfauthor={navstokgap research working notes}}
\newcommand{\tightlist}{\setlength{\itemsep}{0pt}\setlength{\parskip}{0pt}}
\setlength{\emergencystretch}{3em}
\setcounter{secnumdepth}{0}
\begin{document}
\hypertarget{the-confinement-scale-in-three-bands-with-published-numbers-and-the-verification-restated-gauge-invariantly}{%
\section{The confinement scale in three bands, with published numbers,
and the verification restated
gauge-invariantly}\label{the-confinement-scale-in-three-bands-with-published-numbers-and-the-verification-restated-gauge-invariantly}}

Two corrections to the finite-verification statement of
\href{dobrushin-uniqueness-wilson.md}{the Dobrushin note} §4 and
\href{intermediate-region-finite-verification.md}{the
finite-verification note}, and a map of the region with numbers taken
from the literature.

\textbf{Elitzur empties the single-link form of the block criterion.}
Let \(V\) be a box of links and \(s\) a site all of whose links lie in
\(V\). The gauge transformation at \(s\) is a measure-preserving
bijection of the Gibbs measure \(\mu_V^\omega\) for every boundary
condition \(\omega\), so the marginal of any link with an endpoint at
\(s\) is invariant under all left (or right) translations of \(SU(3)\)
and is therefore the Haar measure, whatever \(\omega\) is. The
single-link influence \(\rho_V(x,y)\) vanishes for every link \(x\) with
an interior endpoint, and the sum
\(\sum_{x\in V}\sum_{y\notin V}\rho_V(x,y)\) receives contributions only
from links with both endpoints on the boundary layer. That form of the
criterion yields uniqueness and carries no information about the
interior; the verification must use the \textbf{mixing form on
sub-blocks}, which is the form equivalent to complete analyticity and
the one that yields exponential decay: for every \(W\subseteq V\) and
every \(y\notin V\),
\[\big\|\mu_V^{\omega}\big|_W-\mu_V^{\omega'}\big|_W\big\|_{\rm TV}\ \le\ C\,|W|\,e^{-\gamma\,d(W,y)}
\qquad(\omega,\omega'\text{ differing at }y).\] Since the gauge-variant
part of the marginal on \(W\) is uniform over the gauge orbit at the
interior sites of \(W\), the content of this condition is the
conditional distribution of the gauge-invariant functions of the links
in \(W\), the traces of small Wilson loops, and that is how a
verification should be organized. The single-link Dobrushin condition of
\href{dobrushin-uniqueness-wilson.md}{the Dobrushin note} §2 is
unaffected, because it conditions on all other links and no gauge
symmetry survives the conditioning.

\textbf{The boxes are small where the verification is needed.} With the
scale \(r_0/a\) of Necco and Sommer (Nucl. Phys. B 622 (2002) 328) and
the scalar glueball mass \(r_0m_{0^{++}}=4.21\) of Morningstar and
Peardon (Phys. Rev.~D 60 (1999) 034509), both at metadata level and used
for orientation only, the glueball correlation length in lattice units
is

\begin{longtable}[]{@{}
  >{\raggedright\arraybackslash}p{(\columnwidth - 10\tabcolsep) * \real{0.1667}}
  >{\raggedright\arraybackslash}p{(\columnwidth - 10\tabcolsep) * \real{0.1667}}
  >{\raggedright\arraybackslash}p{(\columnwidth - 10\tabcolsep) * \real{0.1667}}
  >{\raggedright\arraybackslash}p{(\columnwidth - 10\tabcolsep) * \real{0.1667}}
  >{\raggedright\arraybackslash}p{(\columnwidth - 10\tabcolsep) * \real{0.1667}}
  >{\raggedright\arraybackslash}p{(\columnwidth - 10\tabcolsep) * \real{0.1667}}@{}}
\toprule\noalign{}
\begin{minipage}[b]{\linewidth}\raggedright
\(\beta_W\)
\end{minipage} & \begin{minipage}[b]{\linewidth}\raggedright
\(g^2=6/\beta_W\)
\end{minipage} & \begin{minipage}[b]{\linewidth}\raggedright
\(r_0/a\)
\end{minipage} & \begin{minipage}[b]{\linewidth}\raggedright
\(m_{0^{++}}a\)
\end{minipage} & \begin{minipage}[b]{\linewidth}\raggedright
\(\xi/a\)
\end{minipage} & \begin{minipage}[b]{\linewidth}\raggedright
\(\sigma a^2\)
\end{minipage} \\
\midrule\noalign{}
\endhead
\bottomrule\noalign{}
\endlastfoot
\(5.7\) & \(1.05\) & \(2.92\) & \(1.44\) & \(0.70\) & \(0.16\) \\
\(6.0\) & \(1.00\) & \(5.37\) & \(0.78\) & \(1.28\) & \(0.047\) \\
\(6.2\) & \(0.97\) & \(7.38\) & \(0.57\) & \(1.75\) & \(0.025\) \\
\(6.4\) & \(0.94\) & \(9.74\) & \(0.43\) & \(2.31\) & \(0.014\) \\
\end{longtable}

using \(r_0^2\sigma\simeq1.35\) for the string tension. At
\(\beta_W=5.7\) the glueball correlation length is below one lattice
spacing, and a mixing verification would concern boxes of side \(3\) to
\(5\), that is \(216\) to \(2000\) links, in place of the \(10^4\) to
\(10^6\) links estimated earlier from a correlation length of \(1\) to
\(10\). The slowest influence is carried by the flux-tube channel,
\(\sigma a^2\simeq0.16\) per lattice unit at \(\beta_W=5.7\), which sets
the box side rather than the glueball mass. Constants explicit where
they are ours; nothing promoted.

\hypertarget{three-bands-of-the-coupling}{%
\subsection{1. Three bands of the
coupling}\label{three-bands-of-the-coupling}}

\begin{longtable}[]{@{}
  >{\raggedright\arraybackslash}p{(\columnwidth - 4\tabcolsep) * \real{0.3333}}
  >{\raggedright\arraybackslash}p{(\columnwidth - 4\tabcolsep) * \real{0.3333}}
  >{\raggedright\arraybackslash}p{(\columnwidth - 4\tabcolsep) * \real{0.3333}}@{}}
\toprule\noalign{}
\begin{minipage}[b]{\linewidth}\raggedright
band
\end{minipage} & \begin{minipage}[b]{\linewidth}\raggedright
\(\beta_W\)
\end{minipage} & \begin{minipage}[b]{\linewidth}\raggedright
status
\end{minipage} \\
\midrule\noalign{}
\endhead
\bottomrule\noalign{}
\endlastfoot
A & \(<0.0135\) & gapped, proved by hand: Dobrushin single-link
condition, \(g^2>444\) \\
B & \(0.0135\) to about \(5.7\) & mixing with \(\xi/a<1\) by the
published data; no rigorous method reaches it \\
C & above about \(5.7\) & \(\xi/a\) grows, asymptotic scaling sets in
near \(6\) to \(6.5\); the renormalization group's domain \\
\end{longtable}

\emph{Band B is a sharpness problem.} The theory there is, by the data,
more strongly mixing than anything band A contains, yet no expansion
converges with rigorous constants: the polymer expansion of
\href{wilson-strong-coupling-explicit.md}{the Wilson note} reaches
\(\beta_W\simeq0.0057\), the Dobrushin condition \(0.0135\), and an
ideal entropy constant for plaquette surfaces, of order \(10\) to \(15\)
in place of \(20e\), would bring the expansion to \(\beta_W\sim0.1\).
The strong-coupling series itself has singularities near
\(\beta_W\simeq3.5\) to \(4\), so even the true radius of convergence
ends inside band B. The gap between \(0.1\) and \(5.7\) is a factor of
\(50\) in \(\beta_W\) with no small parameter, in a regime where the
correlation length is less than one lattice unit. This is the band the
finite verification addresses, and Section 2 says what it costs.

\emph{Band C is the renormalization group's.} From \(g^2=1/2\)
(\(\beta_W=12\)) to \(\beta_W=6\) (\(g^2=1\)) is \(1/0.0966\simeq10\)
one-loop doublings, over which \(\xi/a\) falls from its
asymptotic-scaling value to about one. The stretch where \(\xi/a\)
passes from about \(10\) to about \(1\) is three doublings wide at one
loop. The physically nontrivial part of the problem, the place where the
gap forms, is this stretch; everything in band B is, physically, already
gapped with a gap of order \(\hbar c/a\).

\hypertarget{what-the-verification-costs-restated}{%
\subsection{2. What the verification costs,
restated}\label{what-the-verification-costs-restated}}

A mixing verification at \(\beta_W\simeq5.7\) with boxes of side \(R=3\)
to \(5\) involves rigorous bounds on the conditional distribution of
small Wilson loops inside \(V\), with \(216\) to \(2000\) link
variables, uniformly over the boundary condition, with the bound
decaying in the distance from the boundary change at a rate consistent
with \(\sigma a^2\simeq0.16\). The uniformity over boundary conditions
is a supremum over a compact set of dimension \(8|\partial V|\), of
order \(10^3\) to \(10^4\), of a function defined by an integral over
\(8|V|\) dimensions. No rigorous numerical method known to the author
evaluates such a quantity with certified bounds today; the problem is
nevertheless three to four orders of magnitude smaller in the number of
variables than the earlier estimate, and it sits at a coupling where the
theory's correlation length is below the lattice spacing, so that a
certified bound on a single small box would settle band B at that
coupling.

\hypertarget{b.-the-meeting-point-located-five-order-one-steps-around-beta_wsimeq6}{%
\subsection{\texorpdfstring{2b. The meeting point located: five
order-one steps around
\(\beta_W\simeq6\)}{2b. The meeting point located: five order-one steps around \textbackslash beta\_W\textbackslash simeq6}}\label{b.-the-meeting-point-located-five-order-one-steps-around-beta_wsimeq6}}

The two unknown numbers of
\href{intermediate-region-finite-verification.md}{the
finite-verification note} §4 can be placed on the map. Read the three
bands as blocking steps, one doubling of the lattice spacing each.

\emph{Band C, from the weak side.} From \(g^2=1/2\) to
\(\beta_W\simeq6\) is ten one-loop doublings, and only the last three,
in which \(\xi/a\) passes from about \(10\) to about \(1\), are outside
the reach of a small-field expansion in principle; the first seven are
the regime of the constructive programme, unavailable with constants but
perturbative in kind.

\emph{Band B, crossed by three steps.} The coarse plaquette at scale
\(2^ka\) is a \(2^k\times2^k\) Wilson loop of the fine lattice, whose
expectation by the area law with \(\sigma a^2\simeq0.16\) at
\(\beta_W=5.7\) is about \(e^{-0.16\cdot4^k}\): \(0.53\), \(0.077\),
\(3.6\times10^{-5}\) for \(k=1,2,3\). Band A needs an activity below
about \(2\times10^{-3}\), so three doublings cross band B. The
Migdal--Kadanoff recursion (Kadanoff, Ann. Phys. 100 (1976) 359;
metadata level, a heuristic in four dimensions), which sends the
activity to its fourth power per doubling, gives two and overestimates
the decrease; the area-law count is the one to use. Rigorously, a
block-spin step is well defined and its image is Gibbsian with a
summable interaction in the strong-mixing regime (van Enter, Fernández
and Sokal, J. Stat. Phys. 72 (1993) 879; metadata level), so once mixing
is known at the top of band B the three steps that cross it are
standard, and band A is proved by hand.

\emph{The meeting point.} The verification is therefore needed at
\textbf{one coupling}, \(\beta_W\simeq5.7\) to \(6\), that is
\(g^2\simeq1\), where \(\xi/a\simeq1\), on boxes of side \(3\) to \(5\);
its openness radius is the tolerance the renormalization group must
meet, and everything below it follows. In the notation of the
finite-verification note, \(g_{\rm DS}^2\simeq1\) conditional on that
single certified box, and the open requirement on the other side is
\(g_{\rm RG}^2\ge1\): the small-field renormalization must be carried
three doublings past the point where \(\xi/a\simeq10\).

\textbf{The whole problem in one sentence.} The mass gap for \(SU(3)\)
is the control of about six consecutive blocking steps of the lattice
theory at order-one coupling, uniformly in the volume: three from the
weak side, where \(\xi/a\) falls from \(10\) to \(1\), and three from
the strong side, which are standard once mixing is certified at the
coupling in between. \href{mass-gap-conditional-theorem.md}{The
conditional theorem} states this as two hypotheses and an explicit
conclusion. Everything else on the map is either perturbative in kind or
proved with explicit constants.

\hypertarget{what-this-changes}{%
\subsection{3. What this changes}\label{what-this-changes}}

The finite verification is restated gauge-invariantly and at its true
size. The decision it asks for is unchanged in kind and smaller in
scale. Bands A and C are as before; band B is where a certified
computation on a box of a few hundred links would connect the
hand-proved region to the physical crossover, and band C remains the
renormalization group's, ten doublings from \(g^2=1/2\) to the crossover
at one loop. The glueball gap along the trajectory is then, in physical
units, \(m_{0^{++}}=4.21/r_0\simeq1.66\,\hbar c\,{\rm fm}^{-1}\) by the
same data, which is what a proof would have to reproduce as a positive
number.

\hypertarget{consequence-for-state}{%
\subsection{4. Consequence for STATE}\label{consequence-for-state}}

The Dobrushin note's specification is corrected: the block criterion in
its single-link form is emptied by Elitzur and the mixing form on
sub-blocks, organized around the conditional distribution of small
Wilson loops, is the one to verify. The verification's size is corrected
downward to boxes of side \(3\) to \(5\) at \(\beta_W\simeq5.7\), where
\(\xi/a<1\) by the published data. The three bands separate the
sharpness problem (B, a factor \(50\) in \(\beta_W\) with \(\xi/a<1\))
from the physical problem (C, three doublings where \(\xi/a\) passes
from \(10\) to \(1\)).
\end{document}
